\section{Benchmarks}
\label{sec:benchmarks}

\subsection{Method}

Every number in this section comes from one run of the benchmark driver shipped with the
repository, \code{benchmarks/run\_benchmarks.sh}, on the machine in Table~\ref{tab:hardware}.
Each point is the median of three timed blocks of steady-state calls after a warm-up, with the
device synchronised once per block. The micro-benchmarks use synthetic Boolean data at 20\%
density, resident on the device unless stated otherwise. The driver writes a JSON record per
measurement and regenerates the figures and tables from it, so the whole section can be
reproduced --- or replaced with another machine's numbers --- with one command.

\begin{table}[htbp]
\centering
\small\sffamily\sloppy
\caption{Measurement environment.}
\label{tab:hardware}
\begin{tabularx}{0.86\linewidth}{@{}L{3.2cm}>{\raggedright\arraybackslash}X@{}}
\ttoprule
\thead{CPU} & AMD Ryzen 9 3950X, 16 cores / 32 threads (PyTorch using 16 threads)\\
\thead{GPU} & NVIDIA GeForce RTX 3090 (24\,GB)\\
\thead{Software} & PyTorch 2.10.0+cu128 / CUDA 12.8, Python 3.10.11, \ttlibname 0.1.2\\
\thead{Reference config} & 784 Boolean features, 500 clauses per class, 10 classes, unless
  stated otherwise\\
\ttbotrule
\end{tabularx}
\end{table}

\begin{ttcallout}[What this section does and does not measure]
It measures \emph{this} library on two devices: the same Python code, the same algorithm, one
run on a CPU and one on a GPU. It is not a comparison against other Tsetlin machine
implementations --- no such comparison was run --- and it is not a state-of-the-art accuracy
claim. Where accuracy appears it is there to show that the batched update learns what it
should, not to compete with 400-epoch results from the literature.
\end{ttcallout}

\subsection{Throughput against batch size}

One \code{update()} call touches every automaton in the machine, so its cost depends far more
on the \emph{size of the machine} than on how many examples are in the batch. That is the
whole point of the reformulation in Section~\ref{sec:batched}, and
Figure~\ref{fig:bench-batch} is what it buys: on the GPU, training throughput rises from 435
examples/s at batch 1 to 219{,}850 at batch 4096 --- a factor of 505 for feeding the same
kernel more work.

\begin{figure}[htbp]
\centering
% Figure: throughput vs batch size, training and inference, CPU vs GPU (MNIST-scale machine).
\pgfplotstableread[col sep=comma]{data/batch-mnist.csv}\batchmnist
\pgfplotstableread[col sep=comma]{data/batch-xor.csv}\batchxor

\begin{tikzpicture}
\begin{groupplot}[
    ttaxis,
    group style={group size=2 by 1, horizontal sep=1.45cm},
    width=0.5\linewidth, height=5.6cm,
    xmode=log, ymode=log,
    log basis x=2,
    xlabel={batch size},
    xtick={1,4,16,64,256,1024,4096},
    xticklabels={1,4,16,64,256,1k,4k},
    legend pos=north west,
  ]

  % ---------------------------------------------------------------- training
  \nextgroupplot[title={Training}, ylabel={examples / s}]
    \addplot[ttgpu] table[x=batch, y=gpu_train] {\batchmnist};
    \addlegendentry{GPU (RTX 3090)}
    \addplot[ttcpu] table[x=batch, y=cpu_train] {\batchmnist};
    \addlegendentry{CPU (16 threads)}
    \addplot[ttthird, dashed, mark=none] table[x=batch, y=gpu_train] {\batchxor};
    \addlegendentry{tiny machine, GPU}
    \addplot[color=ttgray, dotted, line width=1.1pt, mark=none]
      table[x=batch, y=cpu_train] {\batchxor};
    \addlegendentry{tiny machine, CPU}

  % ---------------------------------------------------------------- inference
  \nextgroupplot[title={Inference}, ylabel={examples / s}]
    \addplot[ttgpu] table[x=batch, y=gpu_infer] {\batchmnist};
    \addlegendentry{GPU (RTX 3090)}
    \addplot[ttcpu] table[x=batch, y=cpu_infer] {\batchmnist};
    \addlegendentry{CPU (16 threads)}
\end{groupplot}
\end{tikzpicture}

\caption{Throughput against batch size for a 784-feature, 500-clause-per-class machine
  (log--log). Training on the GPU scales almost linearly with batch size until the device
  saturates. The dashed line is the counter-example: a 12-feature, 10-clause machine is less
  work than a kernel launch, and only overtakes the CPU past batch 1024.}
\label{fig:bench-batch}
\end{figure}

At batch 256 the same model trains at 88{,}908 examples/s on the GPU against 408 on the CPU
--- $218\times$. Inference behaves differently: the CPU forward pass is already cheap, so the
gap there tops out around $28\times$, reaching $1.05\times10^{6}$ examples/s on the GPU.

The dashed series is the honest counterpoint. A Noisy-XOR-scale machine (12 features, 10
clauses per class) is \emph{slower} on the GPU than on the CPU for every batch size below
1024. Small machines belong on a CPU, and \ttlibname makes that a one-word change rather than
a different package.

\subsection{Throughput against model size}

Both axes of model size cost the same thing --- an update is $O(C \cdot 2F)$ element-wise work
--- so both throughput curves fall roughly as $1/x$. The interesting quantity is the
\emph{ratio}, in Figure~\ref{fig:bench-scale}: the CPU falls off immediately while the GPU
stays flat until the machine is large enough to fill it, so the advantage \emph{widens} with
model size, from $51\times$ at 50 clauses per class to $273\times$ at 4000, and from
$71\times$ at 128 features to $286\times$ at 8192.

\begin{figure}[htbp]
\centering
% Figure: how the GPU advantage grows with model size (clause budget and feature count).
\pgfplotstableread[col sep=comma]{data/clauses.csv}\clausetbl
\pgfplotstableread[col sep=comma]{data/features.csv}\featuretbl

\begin{tikzpicture}
\begin{groupplot}[
    ttaxis,
    group style={group size=2 by 1, horizontal sep=1.45cm},
    width=0.5\linewidth, height=5.6cm,
    xmode=log, log basis x=2,
    ymin=0, ymax=320,
    ylabel={GPU speedup over CPU ($\times$)},
    legend pos=north west,
  ]

  \nextgroupplot[title={Clause budget}, xlabel={clauses per class},
                 xtick={50,100,250,500,1000,2000,4000},
                 xticklabels={50,100,250,500,1k,2k,4k}]
    \addplot[ttgpu] table[x=clauses, y=train_speedup] {\clausetbl};
    \addlegendentry{training}
    \addplot[color=ttdeep, dashed, mark=triangle*, mark size=2.2pt, line width=1.1pt,
             mark options={fill=ttamber, draw=ttdeep, solid}]
      table[x=clauses, y=infer_speedup] {\clausetbl};
    \addlegendentry{inference}

  \nextgroupplot[title={Feature count}, xlabel={Boolean features},
                 xtick={128,256,512,1024,2048,4096,8192},
                 xticklabels={128,256,512,1k,2k,4k,8k}]
    \addplot[ttgpu] table[x=features, y=train_speedup] {\featuretbl};
    \addlegendentry{training}
    \addplot[color=ttdeep, dashed, mark=triangle*, mark size=2.2pt, line width=1.1pt,
             mark options={fill=ttamber, draw=ttdeep, solid}]
      table[x=features, y=infer_speedup] {\featuretbl};
    \addlegendentry{inference}
\end{groupplot}
\end{tikzpicture}

\caption{The GPU advantage grows with the model. Big machines --- the ones the Tsetlin machine
  literature actually uses for images --- are exactly the ones worth moving to a GPU.}
\label{fig:bench-scale}
\end{figure}

This matters for a practical reason. The published configurations that reach competitive
accuracy on images use thousands of clauses per class; on a CPU those are overnight runs, and
the resulting iteration speed is what keeps people away. At 2000 clauses per class the GPU
trains at 25{,}934 examples/s against 96 on the CPU.

\subsection{The model zoo}

Weighted and coalesced variants cost about what the plain machine costs --- they change the
output layer and the feedback policy, not the per-clause work that dominates
(Table~\ref{tab:zoo}). The convolutional models are the outlier at roughly $4\times$ the cost
per example, which is arithmetic rather than overhead: a $10\times10$ window over a
$28\times28$ image is 361 patches, and every patch is evaluated against every clause.

\begin{table}[htbp]
\centering
\small\sffamily\sloppy
\caption{All models, 784 features, batch 128. Throughput in examples/s.}
\label{tab:zoo}
\begin{tabularx}{\linewidth}{@{}>{\raggedright\arraybackslash}X R{1.5cm}R{1.6cm}R{1.3cm}R{1.6cm}R{1.7cm}R{1.3cm}@{}}
\ttoprule
& \multicolumn{3}{c}{\thead{Training}} & \multicolumn{3}{c}{\thead{Inference}}\\
\thead{Model} & \thead{CPU} & \thead{GPU} & \thead{speedup} & \thead{CPU} & \thead{GPU} & \thead{speedup}\\
\ttmidrule
\clsname{TsetlinMachine}            & 268 & 49{,}638 & 185$\times$ & 30{,}719 & 526{,}882 & 17$\times$\\
\clsname{TsetlinMachine} (weighted) & 268 & 49{,}475 & 185$\times$ & 32{,}227 & 491{,}498 & 15$\times$\\
\clsname{CoalescedTsetlinMachine}   & 178 & 43{,}579 & 244$\times$ & 34{,}222 & 526{,}083 & 15$\times$\\
\clsname{RegressionTsetlinMachine}  & 214 & 56{,}822 & 266$\times$ & 44{,}567 & 577{,}994 & 13$\times$\\
\clsname{ConvTsetlinMachine}        & 137 & 13{,}059 &  95$\times$ &      140 &  17{,}942 & 128$\times$\\
\clsname{ConvCoalescedTsetlinMachine} & 132 & 13{,}137 & 100$\times$ & 138 & 17{,}826 & 129$\times$\\
\ttbotrule
\end{tabularx}
\end{table}

GPU peak memory is 262\,MB for the flat models and 298\,MB for the convolutional ones at this
configuration --- small enough that clause budget, not memory, is the binding constraint on
most cards.

\subsection{Where an update actually goes}

Figure~\ref{fig:bench-breakdown} breaks one \code{update()} call into stages. The result is
unambiguous: \code{apply\_feedback} --- the stochastic increment and decrement of the whole
$(C,2F)$ state --- is 92\% of the 623\,ms CPU call and 42\% of the 2.8\,ms GPU call. A
standalone \code{torch.binomial} draw of the same shape costs 266.8\,ms on the CPU against
0.22\,ms on the GPU, making it the single most expensive operation in the commit.

\begin{figure}[htbp]
\centering
% Figure: where one update() call spends its time, CPU vs GPU.
% A dot plot rather than bars: on a log axis a bar's length depends on the chosen origin,
% which would be misleading here. The connecting line is the CPU-GPU gap.
\begin{tikzpicture}
\begin{axis}[
    ttaxis,
    width=0.74\linewidth, height=6.6cm,
    xmode=log,
    xmin=0.015, xmax=4000,
    xtick={0.1,1,10,100,1000},
    xticklabels={0.1,1,10,100,1000},
    xlabel={milliseconds per \code{update()} call (log scale)},
    symbolic y coords={{apply feedback}, {accumulator alloc}, {feedback counts},
                       {refresh include}, {votes + select}, evaluate, encode},
    ytick=data,
    enlarge y limits=0.13,
    ymajorgrids=false,
    legend style={at={(0.5,1.14)}, anchor=north, legend columns=2},
  ]

  % CPU-to-GPU gap for each stage
  \draw[ttgray!35, line width=1.1pt] (axis cs:0.035,encode)               -- (axis cs:0.078,encode);
  \draw[ttgray!35, line width=1.1pt] (axis cs:0.207,evaluate)             -- (axis cs:5.183,evaluate);
  \draw[ttgray!35, line width=1.1pt] (axis cs:0.474,{votes + select})     -- (axis cs:7.437,{votes + select});
  \draw[ttgray!35, line width=1.1pt] (axis cs:0.266,{refresh include})    -- (axis cs:14.493,{refresh include});
  \draw[ttgray!35, line width=1.1pt] (axis cs:0.660,{feedback counts})    -- (axis cs:14.229,{feedback counts});
  \draw[ttgray!35, line width=1.1pt] (axis cs:0.113,{accumulator alloc})  -- (axis cs:15.490,{accumulator alloc});
  \draw[ttgray!35, line width=1.1pt] (axis cs:1.166,{apply feedback})     -- (axis cs:575.332,{apply feedback});

  \addplot[only marks, mark=square*, mark size=2.6pt, draw=ttslate, fill=ttslatel,
           point meta=explicit symbolic, nodes near coords,
           nodes near coords style={font=\sffamily\tiny, color=ttslate, anchor=west, xshift=3pt}]
    coordinates {
      (0.078,encode) [0.08] (5.183,evaluate) [5.18] (7.437,{votes + select}) [7.44]
      (14.493,{refresh include}) [14.49] (14.229,{feedback counts}) [14.23]
      (15.490,{accumulator alloc}) [15.49] (575.332,{apply feedback}) [575.3]};
  \addlegendentry{CPU, 16 threads}

  \addplot[only marks, mark=*, mark size=2.6pt, draw=ttdeep, fill=ttorange,
           point meta=explicit symbolic, nodes near coords,
           nodes near coords style={font=\sffamily\tiny, color=ttdeep, anchor=east, xshift=-3pt}]
    coordinates {
      (0.035,encode) [0.04] (0.207,evaluate) [0.21] (0.474,{votes + select}) [0.47]
      (0.266,{refresh include}) [0.27] (0.660,{feedback counts}) [0.66]
      (0.113,{accumulator alloc}) [0.11] (1.166,{apply feedback}) [1.17]};
  \addlegendentry{GPU, RTX 3090}
\end{axis}
\end{tikzpicture}

\caption{Stage timings for one update (5000 clauses, 784 features, batch 256). Note the log
  scale: the CPU total is 623.4\,ms and the GPU total 2.8\,ms. Everything before the commit
  --- encoding, clause evaluation, votes and the feedback-count matmuls --- is comparatively
  free.}
\label{fig:bench-breakdown}
\end{figure}

This is also the quantitative argument for batching. A batched update draws once per commit;
a sequential one draws once per example. Whatever the batch size, the expensive part of the
call happens exactly once.

\subsection{CPU threads}

More CPU threads buy surprisingly little in training --- $1.4\times$ from 1 to 32 threads
(Figure~\ref{fig:bench-threads}). The batched update is a chain of large element-wise
operations over the state, and the binomial draw that dominates it does not thread.
Inference scales much better, $9.6\times$ over the same range, because the forward pass is
one large matmul-shaped reduction. If a CPU run is too slow, the lever is the clause budget
or the batch size, not \code{torch.set\_num\_threads}.

\begin{figure}[htbp]
\centering
% Figure: CPU thread scaling — training barely benefits, inference does.
\begin{tikzpicture}
\begin{axis}[
    ttaxis,
    width=0.56\linewidth, height=5.0cm,
    xmode=log, log basis x=2, ymode=log,
    xlabel={\code{torch} threads}, ylabel={examples / s},
    xtick={1,2,4,8,16,32}, xticklabels={1,2,4,8,16,32},
    legend pos=north west,
  ]
  \pgfplotstableread[col sep=comma]{data/threads.csv}\threadtbl
  \addplot[ttcpu] table[x=threads, y=infer] {\threadtbl};
  \addlegendentry{inference ($9.6\times$ over the range)}
  \addplot[ttgpu, mark=triangle*] table[x=threads, y=train] {\threadtbl};
  \addlegendentry{training ($1.4\times$)}
\end{axis}
\end{tikzpicture}

\caption{CPU thread scaling (5000 clauses, 784 features, batch 256).}
\label{fig:bench-threads}
\end{figure}

\subsection{Where the data lives}

Keeping the Boolean dataset in host memory and copying each mini-batch across costs at most
2\% of throughput on this machine --- the update is long enough to hide the transfer. Streaming
a dataset that does not fit in VRAM is therefore free in practice. The copy that \emph{does}
bite is the one in the other direction, into a CPU-resident encoder, which is the hazard
described in Section~\ref{sec:design}.

\subsection{End to end}

Micro-benchmarks flatter whichever device wins the micro-benchmark, so Table~\ref{tab:e2e}
reports whole training runs with \code{Trainer.fit} included --- data loading, epoch loop,
validation and all.

\begin{table}[htbp]
\centering
\small\sffamily\sloppy
\caption{End-to-end MNIST runs. Accuracy matches across devices to within run-to-run noise,
  as it should: only the wall clock moves.}
\label{tab:e2e}
\begin{tabularx}{\linewidth}{@{}>{\raggedright\arraybackslash}X L{1.5cm}R{1.9cm}R{1.5cm}R{2.0cm}R{1.7cm}@{}}
\ttoprule
\thead{Model} & \thead{Device} & \thead{train ex.} & \thead{epochs} & \thead{s / epoch} & \thead{test acc.}\\
\ttmidrule
\multirow{2}{*}{\clsname{TsetlinMachine} 500/class}
  & CPU  & 60{,}000 & 3 & 114.35 & 0.9542\\
  & GPU  & 60{,}000 & 3 & \textbf{1.52} & 0.9548\\
\ttmidrule
\multirow{2}{*}{\clsname{ConvTsetlinMachine} 200/class, $10\times10$}
  & CPU  & 10{,}000 & 1 & 70.19 & 0.8270\\
  & GPU  & 10{,}000 & 1 & \textbf{0.82} & 0.8240\\
\ttbotrule
\end{tabularx}
\end{table}

A 114-second epoch becomes a 1.5-second one ($75\times$); the convolutional run goes from 70
seconds to 0.8 ($86\times$). That is the difference between a hyper-parameter sweep you plan
and one you simply run.

On accuracy: the flat 500-clause machine reaches 95.8\% after five epochs, and the
convolutional model 97.2\% after three epochs at batch 32 (200 clauses per class, $T=100$,
$s=10$, weighted). The published figures for much larger budgets --- around 98\% for 2000
clauses per class after 400 epochs, and around 99\% for convolutional machines with 8000
clauses per class --- are what the same code should reach given the same budget and patience;
they are not claimed here.

\subsection{What batching costs in accuracy}
\label{sec:fidelity}

The batched update is an approximation, so the honest question is how much fidelity it gives
up. Figure~\ref{fig:bench-fidelity} answers it on Noisy XOR --- the benchmark of the original
paper, deliberately chosen because its generating rule is known, so a failure to learn is
unmistakable. Twenty clauses, $T=15$, $s=3.9$, 100 states, 2000 training examples with 40\%
label noise, three seeds per configuration, with the number of epochs raised as the batch grows
so that every configuration sees the data a comparable number of times.

\begin{figure}[htbp]
\centering
% Figure: what batching costs in accuracy — Noisy XOR, 3 seeds per batch size.
% A strip plot rather than bars: the y axis does not start at zero, and bar lengths on a
% truncated axis would overstate the differences.
\begin{tikzpicture}
\begin{axis}[
    ttaxis,
    width=0.66\linewidth, height=5.6cm,
    symbolic x coords={1, 10, 50, 200},
    xtick=data,
    xlabel={batch size \ (1 = the exact sequential algorithm)},
    ylabel={test accuracy},
    ymin=0.73, ymax=1.03,
    ytick={0.75,0.80,0.85,0.90,0.95,1.00},
    yticklabel style={/pgf/number format/fixed, /pgf/number format/zerofill,
                      /pgf/number format/precision=2},
    enlarge x limits=0.18,
    legend style={at={(0.5,-0.32)}, anchor=north, legend columns=3},
  ]

  % seed-to-seed range
  \draw[ttgray!45, line width=1.2pt] (axis cs:1,0.954)   -- (axis cs:1,0.993);
  \draw[ttgray!45, line width=1.2pt] (axis cs:10,0.931)  -- (axis cs:10,0.993);
  \draw[ttgray!45, line width=1.2pt] (axis cs:50,0.877)  -- (axis cs:50,0.993);
  \draw[ttgray!45, line width=1.2pt] (axis cs:200,0.771) -- (axis cs:200,0.889);

  \addplot[only marks, mark=o, mark size=2.6pt, draw=ttslate, line width=0.9pt]
    coordinates {(1,0.954) (1,0.955) (1,0.993)
                 (10,0.942) (10,0.993) (10,0.931)
                 (50,0.993) (50,0.965) (50,0.877)
                 (200,0.788) (200,0.771) (200,0.889)};
  \addlegendentry{individual seeds}

  \addplot[only marks, mark=diamond*, mark size=3.6pt, fill=ttorange, draw=ttdeep,
           line width=0.7pt, point meta=explicit symbolic, nodes near coords,
           nodes near coords style={font=\sffamily\scriptsize, color=ttdeep,
                                    anchor=west, xshift=5pt}]
    coordinates {(10,0.955) [0.955] (50,0.945) [0.945] (200,0.816) [0.816]};
  \addlegendentry{mean of 3 seeds}

  % batch 1 is labelled below-right so the label clears the reference line it defines
  \addplot[only marks, mark=diamond*, mark size=3.6pt, fill=ttorange, draw=ttdeep,
           line width=0.7pt, point meta=explicit symbolic, nodes near coords, forget plot,
           nodes near coords style={font=\sffamily\scriptsize, color=ttdeep,
                                    anchor=north west, xshift=3pt, yshift=-2pt}]
    coordinates {(1,0.967) [0.967]};

  \addplot[ttamber, dashed, line width=1.0pt, mark=none, forget plot]
    coordinates {(1,0.967) (200,0.967)};
  \addlegendimage{ttamber, dashed, line width=1.0pt, mark=none}
  \addlegendentry{sequential mean}
\end{axis}
\end{tikzpicture}

\caption{Noisy XOR accuracy against batch size (3 seeds). Batches up to about 50 land in the
  same band as the exact sequential algorithm; 200 degrades clearly. Seed-to-seed spread is
  wide on this task --- which is itself the reason to read the individual points rather than
  only the means.}
\label{fig:bench-fidelity}
\end{figure}

Batch sizes of roughly 8--64 keep the accuracy of the sequential algorithm while running one
to two orders of magnitude faster on a GPU. At 200 the degradation is unambiguous. The
practical recommendation the library documents follows directly: start at 16--64 for flat
models and 8--32 for convolutional ones, and if accuracy lags, halve the batch or add epochs
before touching $T$ and $s$.

The cost of exactness is measurable too. At batch 32, \code{feedback\_mode="sequential"} runs
7--28$\times$ slower than the batched update depending on model scale and device, because it
performs 32 commits where the batched path performs one. That is the trade, stated in numbers:
\emph{batching buys one to two orders of magnitude and costs accuracy only once the batch
grows past a few tens of examples.}

\begin{ttkey}
Reformulating the update as event counting plus one binomial draw makes the cost of an update
nearly independent of batch size --- which is what turns a $114$-second MNIST epoch into a
$1.5$-second one, and what makes a clause budget of thousands a routine choice rather than an
overnight one.
\end{ttkey}
